\documentclass[12pt,a4paper]{article}

% Packages
\usepackage[margin=1in]{geometry}
\usepackage{hyperref}
\usepackage{graphicx}
\usepackage{booktabs}
\usepackage{array}
\usepackage{longtable}
\usepackage{enumitem}
\usepackage{titlesec}
\usepackage{fancyhdr}
\usepackage{xcolor}
\usepackage{mdframed}
\usepackage{setspace}
\usepackage{parskip}

% Colors
\definecolor{nuclearblue}{RGB}{0,51,102}
\definecolor{warningred}{RGB}{180,30,30}
\definecolor{lightgray}{RGB}{240,240,240}

% Header/Footer
\setlength{\headheight}{15pt}
\pagestyle{fancy}
\fancyhf{}
\rhead{\small Plant Vogtle Nuclear Power Plant --- Case Study}
\lhead{\small \thepage}
\renewcommand{\headrulewidth}{0.4pt}

% Section formatting
\titleformat{\section}{\large\bfseries\color{nuclearblue}}{}{0em}{}[\titlerule]
\titleformat{\subsection}{\normalsize\bfseries\color{nuclearblue}}{}{0em}{}

% Hyperref setup
\hypersetup{
    colorlinks=true,
    linkcolor=nuclearblue,
    urlcolor=nuclearblue,
    pdftitle={Plant Vogtle Nuclear Power Plant: A Comprehensive Case Study},
    pdfauthor={Research Analysis Report},
}

\begin{document}

% ─────────────────────────── TITLE PAGE ────────────────────────────────────
\begin{titlepage}
\centering
\vspace*{2cm}

{\Huge\bfseries\color{nuclearblue}
Plant Vogtle Nuclear Power Plant\\[0.5em]
A Comprehensive Case Study}

\vspace{1em}
{\large\color{gray} America's Largest Nuclear Power Station ---\\
Cost Overruns, Schedule Delays, and Lessons Learned}

\vspace{2em}
\rule{0.8\textwidth}{1pt}

\vspace{2em}
{\large\bfseries Based on the Following Key Studies:}

\vspace{1em}
\begin{mdframed}[backgroundcolor=lightgray, linecolor=nuclearblue, linewidth=1pt]
\begin{enumerate}[leftmargin=1.5em, label=\textbf{\arabic*.}]
    \item \textbf{``The Truth About Vogtle''} \\
          \textit{Patty Durand, Kim Scott, Glenn Carroll} \\
          Sponsored by: Cool Planet Solutions, Georgia WAND, Nuclear Watch South, \\
          GCV Education Fund, Center for a Sustainable Coast, \\
          Concerned Ratepayers of Georgia (2023--2024)

    \item \textbf{``Ratepayer Robbery --- The True Cost of Plant Vogtle''} \\
          \textit{Georgia Conservation Voters \& Environment Georgia} \\
          December 2021

    \item \textbf{``Southern Company's Troubled Vogtle Nuclear Project''} \\
          \textit{David Schlissel, Director of Resource Planning Analysis} \\
          Institute for Energy Economics and Financial Analysis (IEEFA), January 2022

    \item \textbf{``Vogtle Units 3 \& 4: Twenty-Third Semi-Annual \\
          Construction Monitoring Report''} \\
          \textit{Georgia Power Company / Southern Nuclear Operating Company} \\
          Docket No. 29849, August 2020

    \item \textbf{``Final Report SJR 79: Kentucky Nuclear Energy Development''} \\
          \textit{Kentucky Office of Energy Policy} \\
          Report to Kentucky Legislative Research Commission, November 17, 2023

    \item \textbf{``Commercial Nuclear Power Reactors --- Operating Reactors''} \\
          \textit{U.S. Nuclear Regulatory Commission (NRC)} \\
          Document No. ML25051A096 (Appendix A, NRC Reactor List)
\end{enumerate}
\end{mdframed}

\vfill
{\normalsize Research Analysis Report \\ April 2026}
\end{titlepage}

% ─────────────────────────── TABLE OF CONTENTS ─────────────────────────────
\tableofcontents
\newpage

% ─────────────────────────── 1. INTRODUCTION ───────────────────────────────
\section{Introduction}

Georgia Power's Plant Vogtle, located in Burke County, Georgia near Augusta, is the
largest nuclear power station in the United States. It is the only site in the country
where new commercial nuclear reactors have been constructed in over 30 years,
making it a landmark case study in both nuclear energy policy and large infrastructure
project management.

Plant Vogtle Units 1 and 2 came online in 1987 and 1989, respectively, having
themselves experienced dramatic cost overruns---escalating from an initial estimate
of \$660 million to a final cost of \$8.8 billion. Despite this precedent, Georgia Power
Company, a wholly-owned subsidiary of the Southern Company, sought and received
approval to construct two additional reactors (Units 3 and 4) employing Westinghouse's
new AP1000 passive-safety reactor design.

This case study synthesizes findings from six key reports and documents produced
between 2020 and 2025 by a range of organizations---from advocacy groups and
independent financial analysts to regulatory agencies and utility operators---to provide
a comprehensive account of Plant Vogtle's development, cost, controversy, and
implications for U.S. energy policy.

% ─────────────────────────── 2. PLANT OVERVIEW ─────────────────────────────
\section{Plant Vogtle: Overview and Background}

\subsection{Location and Ownership}

Plant Vogtle is situated on the Savannah River in Burke County, Georgia. The four
owners of Vogtle Units 3 and 4 are:

\begin{itemize}
    \item \textbf{Georgia Power Company} (lead owner and agent) --- a subsidiary of Southern Company
    \item \textbf{Oglethorpe Power Corporation}
    \item \textbf{Municipal Electric Authority of Georgia (MEAG)}
    \item \textbf{City of Dalton, Georgia}
\end{itemize}

Southern Nuclear Operating Company (SNC), acting as agent for Georgia Power,
managed the construction and operations of Units 3 and 4.

\subsection{The AP1000 Reactor Design}

Units 3 and 4 use Westinghouse's AP1000 reactor design, a pressurized water reactor
(PWR) with passive safety systems. When Georgia Power applied to the Georgia Public
Service Commission (PSC) in 2008 for certification, it argued that the AP1000's modular
construction techniques would enable the project to be completed faster and at lower
cost than prior nuclear builds. This claim proved to be significantly overstated.

\subsection{Historical Context: Units 1 and 2}

The construction history of Vogtle Units 1 and 2 provided a clear warning sign
that was ultimately ignored. Initial cost estimates for Units 1 and 2 stood at
\$660 million; the final cost ballooned to \$8.8 billion by the time they came online
in 1987 and 1989---an overrun exceeding 1,200\%. A U.S. Department of Energy
study of 75 reactors built between 1966 and 1977 found average construction costs
ran 207\% over initial estimates. The historical record provided ample evidence that
new nuclear construction was prone to severe cost escalation.

% ─────────────────────────── 3. CONSTRUCTION TIMELINE ─────────────────────
\section{Construction Timeline and Schedule Delays}

\subsection{Original Schedule}

Georgia Power proposed the following commercial operation dates (CODs) when
seeking certification in 2008:

\begin{center}
\begin{tabular}{lcc}
\toprule
\textbf{Unit} & \textbf{Originally Estimated COD} & \textbf{Actual / Revised COD} \\
\midrule
Unit 3 & April 1, 2016 & July 2023 \\
Unit 4 & April 1, 2017 & Late 2023 / Early 2024 \\
\bottomrule
\end{tabular}
\end{center}

Both units entered commercial operation more than \textbf{six years behind schedule}.

\subsection{Key Construction Milestones and Setbacks}

The Twenty-Third Semi-Annual Vogtle Construction Monitoring Report (August 2020,
Docket No. 29849) covered the period January 1 to June 30, 2020 and reported:

\begin{itemize}
    \item Southern Nuclear Operating Company received no NRC Notices of Violation
          during the reporting period and maintained a \textbf{green (favorable) status}
          under the NRC's Construction Reactor Oversight Process (cROP).
    \item The site sustained more than \textbf{34 million man-hours} without a lost-time
          incident before an incident occurred in April 2020.
    \item The \textbf{COVID-19 pandemic} caused significant workforce disruptions in
          2020, including spikes in infections (over 800 workers tested positive), workforce
          reductions, and increased costs and productivity challenges.
    \item Despite pandemic-related setbacks, essential construction activities including
          the Passive Containment Cooling Water Storage Tank (CB20) on the Unit 3
          Shield Building continued.
\end{itemize}

\subsection{Westinghouse Bankruptcy (2017)}

The most dramatic event in Vogtle's construction history was the March 2017
bankruptcy filing by Westinghouse Electric Company, the contractor responsible for
engineering, procurement, and construction. The bankruptcy was triggered in part by
billions of dollars in losses on the Vogtle and V.C. Summer nuclear construction projects.

In South Carolina, an identical AP1000 project (SCANA's V.C. Summer) was abandoned
in July 2017 following nine consecutive rate increases. Investigations revealed that
Westinghouse and utility executives had misrepresented project progress and costs,
ultimately resulting in criminal charges and jail time for two SCANA executives.

In Georgia, the PSC voted to continue the Vogtle project despite the Westinghouse
bankruptcy, without imposing meaningful consumer protections against further overruns.

% ─────────────────────────── 4. COST ANALYSIS ──────────────────────────────
\section{Cost Analysis and Financial Impact}

\subsection{Cost Escalation}

\begin{center}
\begin{tabular}{lcc}
\toprule
\textbf{Year} & \textbf{Estimated Total Cost} & \textbf{Source} \\
\midrule
2009 (original approval) & \$14.0 billion & Georgia PSC Certification \\
2017 (post-bankruptcy) & \$25+ billion & Georgia PSC Proceeding \\
2021 & \$32+ billion & GCV Ratepayer Robbery Report \\
2022 & \$30+ billion & IEEFA / Schlissel Report \\
\bottomrule
\end{tabular}
\end{center}

The total cost of Units 3 and 4 more than \textbf{doubled} from the original estimate
of \$14 billion to above \$30--32 billion, making Plant Vogtle the most expensive
power plant ever built on Earth at the time of its completion.

\subsection{Cost Compared to Alternatives}

The GCV ``Ratepayer Robbery'' report (December 2021) calculated that the cost per
megawatt-hour of energy from Vogtle was \textbf{6 to 9 times higher} than comparable
energy from renewable sources or natural gas that could have produced the same
2,200 MW of generating capacity.

\subsection{The Nuclear Construction Cost Recovery (NCCR) Tariff}

Georgia was one of a small number of states with a ``Construction Work in Progress''
(CWIP) law allowing utilities to collect construction costs from ratepayers \textit{before}
a plant enters commercial operation. The NCCR tariff allowed Georgia Power to
collect revenues from customers throughout the construction period, effectively
shifting financial risk from shareholders to ratepayers. Between 2009 and 2021,
Georgia Power collected multiple NCCR-based rate increases from its customers.

\subsection{Return on Equity and Regulatory Capture}

Critics, including the GCV report, argued that the Georgia PSC allowed Georgia Power
one of the \textbf{highest returns on equity (ROE)} of any utility in the United States.
This structure meant Georgia Power \textit{profited from construction delays}: longer
construction timelines increased the capital base on which the utility earned its
guaranteed rate of return. This created a perverse incentive that contributed to
inadequate oversight of project management.

% ─────────────────────────── 5. REGULATORY OVERSIGHT ──────────────────────
\section{Regulatory Oversight and the Georgia PSC}

\subsection{Role of the Georgia Public Service Commission}

The Georgia Public Service Commission (PSC) is the body responsible for regulating
Georgia Power as a monopoly utility and for protecting ratepayers. The PSC approved
the Vogtle project in March 2009 and voted to continue construction after the
Westinghouse bankruptcy in 2017.

Critics across multiple reports characterized the PSC as suffering from
``\textbf{regulatory capture}''---a phenomenon where a regulatory body acts in the
interests of the industry it regulates rather than the public. Georgia Power itself
referred to this relationship as ``constructive regulation.''

\subsection{Key PSC Decisions}

\begin{itemize}
    \item \textbf{March 2009}: Approved Georgia Power's application to build Units 3 and 4.
    \item \textbf{2017}: Voted to continue construction despite Westinghouse bankruptcy and
          available cancellation option; declined to impose consumer protections.
    \item \textbf{December 2023}: Final PSC vote following completion of the project, amid
          ongoing protests from ratepayer advocacy groups.
\end{itemize}

\subsection{NRC Oversight}

The U.S. Nuclear Regulatory Commission (NRC) monitored construction through its
Construction Reactor Oversight Process (cROP). Per the Twenty-Third VCM Report
(August 2020), Southern Nuclear Operating Company maintained a green (satisfactory)
NRC performance indicator throughout 2020.

The NRC document ML25051A096 lists Plant Vogtle among the United States' licensed
commercial nuclear power reactors, confirming its operational status in the NRC's
national reactor registry.

% ─────────────────────────── 6. STUDY-BY-STUDY ANALYSIS ────────────────────
\section{Study-by-Study Analysis}

\subsection{Study 1: ``The Truth About Vogtle''}

\textbf{Authors:} Patty Durand (Cool Planet Solutions), Kim Scott (Georgia WAND),
Glenn Carroll (Nuclear Watch South)\\
\textbf{Sponsors:} Center for a Sustainable Coast; Concerned Ratepayers of Georgia;
Cool Planet Solutions; GCV Education Fund; Georgia WAND; Nuclear Watch South

This report, the most comprehensive advocacy document in the repository, presents
a critical examination of the entire Vogtle project from initial approval through
completion. Key findings include:

\begin{itemize}
    \item Vogtle 3 and 4 are the \textbf{only new nuclear reactors built in the U.S. in
          over 30 years}.
    \item Cost overruns began immediately after construction started in 2009.
    \item The Westinghouse bankruptcy in March 2017 created a clear decision point
          that the Georgia PSC failed to use to protect ratepayers.
    \item The report systematically debunks seven common myths about nuclear energy:
    \begin{enumerate}
        \item Nuclear energy is clean
        \item Nuclear energy is safe
        \item Nuclear waste is ``no big deal''
        \item Small Modular Reactors (SMRs) are different
        \item Nuclear is required to back up renewables
        \item Nuclear is required to meet future growth
        \item Nuclear is needed to combat climate change
    \end{enumerate}
    \item The report documents \textbf{scandals and litigation} surrounding the project,
          including parallels with the V.C. Summer SCANA scandal in South Carolina
          where executives faced criminal charges.
    \item The U.S. South is described as the ``nuclear hub of the United States,''
          hosting over three dozen operational reactors, weapons manufacturing complexes,
          and nuclear fuel factories.
\end{itemize}

\subsection{Study 2: ``Ratepayer Robbery --- The True Cost of Plant Vogtle''}

\textbf{Authors:} Georgia Conservation Voters \& Environment Georgia\\
\textbf{Date:} December 2021

This report frames the Vogtle project as a financial harm imposed on Georgia
ratepayers. Key findings include:

\begin{itemize}
    \item Plant Vogtle is described as ``\textbf{the most expensive power plant ever
          built on Earth},'' with costs exceeding \$32 billion for 2,200 MW of capacity.
    \item The cost-per-megawatt-hour from Vogtle is 6--9 times higher than
          renewable alternatives.
    \item Georgia Power's system outages \textbf{consistently exceeded U.S. averages}
          for investor-owned utilities, undermining the utility's claim of ``superior service.''
    \item The report details the NCCR tariff mechanism and documents how Georgia
          Power collected rate increases from customers during construction.
    \item The ``Parental Guarantee'' from Southern Company is analyzed, showing
          the limits of shareholder protection for ratepayers.
    \item The report argues that the PSC's 2017 decision to continue construction
          was the \textbf{costliest regulatory failure} in Georgia history.
    \item It concludes that cancellation and replacement with renewables in 2017
          would have been cheaper for customers.
\end{itemize}

\subsection{Study 3: ``Southern Company's Troubled Vogtle Nuclear Project''}

\textbf{Author:} David Schlissel, Director of Resource Planning Analysis, IEEFA\\
\textbf{Date:} January 2022

This report, produced by the Institute for Energy Economics and Financial Analysis
(IEEFA), provides an independent financial and engineering assessment. Key findings:

\begin{itemize}
    \item Total cost had climbed above \$30 billion by January 2022, with completion
          still expected to be \textbf{more than six years behind the original schedule}.
    \item Georgia Power's 2008 application claimed modular AP1000 construction would
          reduce both cost and schedule---a claim that proved false.
    \item The report draws a direct line between the failure of Vogtle 1 \& 2 (which
          cost \$8.8 billion versus an estimate of \$660 million) and the failure of
          Units 3 \& 4, arguing these outcomes were \textbf{predictable} given nuclear
          construction history.
    \item A DOE study of 75 reactors found average cost overruns of 207\%---evidence
          that was available to the PSC before it approved Vogtle 3 \& 4.
    \item Commercial operations for Unit 3 were projected at 2022 and Unit 4 at 2023,
          though actual completion slipped further.
\end{itemize}

\subsection{Study 4: ``Vogtle Units 3 \& 4: Twenty-Third Semi-Annual Construction Monitoring Report''}

\textbf{Issuer:} Georgia Power Company / Southern Nuclear Operating Company\\
\textbf{Date:} August 2020 (covering January--June 2020)\\
\textbf{Docket:} No. 29849

This is an official regulatory filing by the utility itself, submitted to the
Georgia PSC. Key contents:

\begin{itemize}
    \item The report describes the overall safety and quality management culture on site,
          including a commitment to ``\textbf{safety first}.''
    \item As of June 2020, SNC maintained green status under the NRC's cROP and
          had received no Notices of Violation.
    \item COVID-19 significantly impacted the project: over 800 workers tested positive,
          and workforce reductions in April 2020 caused productivity losses, cost increases,
          and milestone delays.
    \item The report shows the Passive Containment Cooling Water Storage Tank (CB20)
          atop the Unit 3 Shield Building being completed---a visible symbol of construction
          progress during a turbulent period.
    \item This report represents the utility's perspective, emphasizing compliance and
          safety culture, in contrast to the critical perspective of the advocacy reports.
\end{itemize}

\subsection{Study 5: ``Final Report SJR 79 --- Kentucky Nuclear Energy Development''}

\textbf{Issuer:} Kentucky Office of Energy Policy\\
\textbf{Date:} November 17, 2023\\
\textbf{Submitted to:} Kentucky Legislative Research Commission

While primarily focused on Kentucky's nuclear energy future, this report provides
important national context drawn in part from the Vogtle experience:

\begin{itemize}
    \item Kentucky currently has \textbf{no operational commercial nuclear reactors},
          having relied on coal and natural gas historically.
    \item The report analyzes the \textbf{value proposition of nuclear energy} in terms
          of reliability, baseload power, and decarbonization.
    \item It examines barriers to nuclear development in Kentucky, including:
    \begin{itemize}
        \item High capital costs and financing risk (informed by Vogtle's experience)
        \item Workforce and supply chain gaps
        \item Regulatory and permitting uncertainty
        \item Public perception challenges
    \end{itemize}
    \item The report recommends creating a \textbf{Kentucky Nuclear Energy Development
          Authority} with near-term goals (0--3 years) and intermediate-term goals (3+ years).
    \item It identifies Small Modular Reactors (SMRs) as a ``gated opportunity''
          for future deployment, contingent on the success of national SMR demonstration
          projects.
    \item The Vogtle project's lessons---particularly the dangers of cost overruns with
          large light-water reactors---are implicitly used to justify a cautious, phased
          approach to nuclear development in Kentucky.
\end{itemize}

\subsection{Study 6: NRC Document ML25051A096 --- ``Commercial Nuclear Power Reactors: Operating Reactors''}

\textbf{Issuer:} U.S. Nuclear Regulatory Commission (NRC)\\
\textbf{Document:} Appendix A, Operating Reactors List

This NRC document places Plant Vogtle in the national context of licensed and
operating commercial nuclear power reactors in the United States. Key information:

\begin{itemize}
    \item The document lists all NRC-licensed operating reactors by plant name,
          licensee, location, docket number, reactor type (e.g., PWR-DRYAMB, BWR),
          NSSS supplier, architect-engineer, constructor, net MWe capacity, and NRC Region.
    \item Plant Vogtle appears as a licensed facility in NRC Region II (Southeast).
    \item The listing confirms that Plant Vogtle Units 3 and 4 are operational, using
          the AP1000 (PWR) design supplied by Westinghouse.
    \item The document contextualizes Vogtle among a national fleet that includes:
          Arkansas Nuclear One (Entergy), Beaver Valley Power Station (Energy Harbor/Vistra),
          Braidwood Station (Constellation), and dozens of others.
    \item The NRC's role as licensing authority---distinct from the state PSC's economic
          regulation---is clarified by this document.
\end{itemize}

% ─────────────────────────── 7. KEY THEMES ─────────────────────────────────
\section{Cross-Cutting Themes and Lessons Learned}

\subsection{Theme 1: The Persistent Pattern of Nuclear Cost Overruns}

Every independent study reviewed here confirms that large-scale nuclear construction
in the United States has been associated with dramatic cost overruns. From Vogtle
Units 1 \& 2 in the 1980s to Units 3 \& 4 in the 2010s--2020s, and across other
projects documented by the DOE and NRC, cost escalation appears to be a structural
feature of nuclear construction rather than an anomaly.

\subsection{Theme 2: Regulatory Capture and Ratepayer Risk}

Multiple reports identify a pattern of regulatory capture at the Georgia PSC. The
NCCR tariff mechanism enabled Georgia Power to pass construction costs to ratepayers
in advance, removing the market discipline that would otherwise incentivize cost
control. The PSC's failure to impose consumer protections following the Westinghouse
bankruptcy in 2017 is highlighted by all three independent studies as a critical
governance failure.

\subsection{Theme 3: The V.C. Summer Comparison}

The abandoned V.C. Summer project in South Carolina serves as a counterfactual
throughout the advocacy reports. Had Georgia followed South Carolina's decision to
cancel the project in 2017, ratepayers would have avoided billions in additional costs.
The criminal accountability for SCANA executives in South Carolina---absent in
Georgia---is cited as evidence of differential regulatory enforcement.

\subsection{Theme 4: Implications for Future Nuclear Development}

The Kentucky SJR 79 report illustrates how Vogtle's troubled history is shaping
nuclear policy discussions in other states. The report's cautious framing of SMRs
as a ``gated opportunity'' reflects lessons drawn from Vogtle: the importance of
proven technology, phased investment, and independent financial oversight before
committing to large capital expenditures.

\subsection{Theme 5: COVID-19 and Infrastructure Resilience}

The VCM Report from August 2020 documents how the COVID-19 pandemic exposed the
vulnerability of large construction projects to workforce disruptions. The virus
caused schedule delays and cost increases at Vogtle, adding to an already troubled
project. This experience points to the need for resilient workforce planning and
contingency budgets in large infrastructure projects.

% ─────────────────────────── 8. DATA SUMMARY TABLE ─────────────────────────
\section{Key Data Summary}

\begin{center}
\begin{longtable}{p{4cm} p{9cm}}
\toprule
\textbf{Parameter} & \textbf{Data} \\
\midrule
\endfirsthead
\toprule
\textbf{Parameter} & \textbf{Data} \\
\midrule
\endhead
\bottomrule
\endfoot
Location & Burke County, Georgia, USA (near Augusta, on the Savannah River) \\[4pt]
Plant Type & Nuclear Power Plant (Pressurized Water Reactor --- AP1000) \\[4pt]
Units 3 \& 4 Capacity & 2,200 MW (combined) \\[4pt]
Owner (Lead) & Georgia Power Company (subsidiary of Southern Company) \\[4pt]
Other Owners & Oglethorpe Power, MEAG, City of Dalton \\[4pt]
Reactor Supplier & Westinghouse Electric Company (AP1000 design) \\[4pt]
NRC Region & Region II (Southeast) \\[4pt]
Original Cost Estimate & \$14 billion (2009 PSC application) \\[4pt]
Final Cost Estimate & \$32+ billion \\[4pt]
Cost Overrun & $\sim$130\% above original estimate \\[4pt]
Original COD (Unit 3) & April 1, 2016 \\[4pt]
Actual COD (Unit 3) & July 2023 \\[4pt]
Schedule Delay & 6+ years behind schedule \\[4pt]
Westinghouse Bankruptcy & March 2017 \\[4pt]
NRC cROP Status (2020) & Green (Satisfactory) \\[4pt]
Man-Hours without LTI & 34+ million (as of April 2020) \\[4pt]
COVID-19 Impact & 800+ positive tests on site; workforce reduction; cost increases \\[4pt]
Cost vs. Renewables & 6--9$\times$ higher per MWh than renewable alternatives \\[4pt]
\end{longtable}
\end{center}

% ─────────────────────────── 9. CONCLUSION ─────────────────────────────────
\section{Conclusion}

Plant Vogtle Units 3 and 4 represent the most significant nuclear energy construction
project in the United States in a generation. As the only new commercial nuclear
reactors built in the country in over 30 years, they have served as both an ambition
and a cautionary tale.

The six studies reviewed in this case study converge on several conclusions:

\begin{enumerate}
    \item \textbf{Cost overruns were predictable.} Historical data from Vogtle 1 \& 2
          and the broader U.S. nuclear fleet clearly indicated that the original \$14 billion
          estimate was unrealistic. The project ultimately cost more than double that amount.

    \item \textbf{Regulatory oversight failed ratepayers.} The Georgia PSC, functioning
          under a model critics describe as regulatory capture, repeatedly made decisions
          that protected utility profits over consumer interests---most critically in 2017
          when the Westinghouse bankruptcy created a legitimate off-ramp.

    \item \textbf{The financial burden fell on ratepayers.} Through the NCCR tariff and
          repeated rate increases, Georgia electricity customers funded the project during
          construction and will continue to pay elevated rates for decades.

    \item \textbf{Alternatives existed and were not chosen.} Multiple analyses demonstrate
          that renewable energy and natural gas alternatives could have provided equivalent
          capacity at a fraction of the cost, especially given the dramatic price reductions
          in solar and wind between 2009 and 2023.

    \item \textbf{Vogtle is now shaping national nuclear policy.} The project's completion,
          despite its cost and schedule failures, has been cited by nuclear advocates as
          proof of concept for AP1000 technology. Other states such as Kentucky are
          examining nuclear options while attempting to learn from Vogtle's mistakes
          through phased approaches and SMR technology.
\end{enumerate}

Plant Vogtle stands as a complex legacy: technically completed but economically
troubled; a proof of advanced reactor technology but also evidence of deep systemic
failures in utility regulation, cost estimation, and infrastructure governance. Its
lessons will shape U.S. energy policy for decades to come.

% ─────────────────────────── REFERENCES ────────────────────────────────────
\section*{References}
\addcontentsline{toc}{section}{References}

\begin{enumerate}[leftmargin=1.5em, label={[\arabic*]}]

    \item Durand, P., Scott, K., \& Carroll, G. (n.d.). \textit{The Truth About Vogtle}.
          Sponsored by Cool Planet Solutions, Georgia WAND, Nuclear Watch South,
          GCV Education Fund, Center for a Sustainable Coast, Concerned Ratepayers
          of Georgia.

    \item Georgia Conservation Voters \& Environment Georgia. (December 2021).
          \textit{Ratepayer Robbery --- The True Cost of Plant Vogtle}.

    \item Schlissel, D. (January 2022). \textit{Southern Company's Troubled Vogtle
          Nuclear Project: Units 3 and 4 Now Expected to Cost More Than \$30 Billion
          and Are at Least Six Years Behind Schedule}. Institute for Energy Economics
          and Financial Analysis (IEEFA).

    \item Georgia Power Company / Southern Nuclear Operating Company. (August 2020).
          \textit{Vogtle Units 3 \& 4: Twenty-Third Semi-Annual Construction Monitoring
          Report}. Docket No. 29849. Submitted to Georgia Public Service Commission.

    \item Kentucky Office of Energy Policy. (November 17, 2023). \textit{Final Report
          SJR 79: Report to the Kentucky Legislative Research Commission Pursuant to
          2023RS SJR 79}.

    \item U.S. Nuclear Regulatory Commission. (2025). \textit{Commercial Nuclear Power
          Reactors: Operating Reactors, Appendix A}. NRC Document ML25051A096.

    \item U.S. Energy Information Administration. (January 1, 1986). \textit{An Analysis
          of Nuclear Power Plant Construction Costs}. Technical Report DOE/EIA-0485.

    \item Georgia Public Service Commission. (August 1, 2008). \textit{Georgia Power
          Company Application for the Certification of Vogtle Generating Units 3 and 4
          and Upgraded Integrated Resource Plan}.

\end{enumerate}

\end{document}
